\subsection{Exception Class Hierarchies and Handler Selection}

Python exceptions are not untyped failure messages. They are objects created from exception classes. This means that exception handling is tied directly to Python's runtime type system. An \texttt{except} clause does not merely catch a string or a numeric error code. It matches an exception object by its class and by that class's position in the exception inheritance hierarchy.

This is why handler order matters. A specific exception class should usually be handled before a more general parent class. If the general handler appears first, it can catch the exception before the more specific handler ever has a chance to run.

\begin{quote}
Exception handling is runtime type selection applied to non-local control flow.
\end{quote}

\subsubsection{Exception Classes as Runtime Types}

An exception object has a type, just like an integer, string, list, tuple, or dictionary has a type.

\begin{verbatim}
try:
    number = int("forty-two")
except ValueError as err:
    print(f"err: {err}")
    print(f"type(err): {type(err)}")
    print(f"is ValueError: {isinstance(err, ValueError)}")
    print(f"is Exception: {isinstance(err, Exception)}")
    print(f"is BaseException: {isinstance(err, BaseException)}")
\end{verbatim}

Possible output:

\begin{verbatim}
err: invalid literal for int() with base 10: 'forty-two'
type(err): <class 'ValueError'>
is ValueError: True
is Exception: True
is BaseException: True
\end{verbatim}

The object bound to \texttt{err} is a \texttt{ValueError} object. Because \texttt{ValueError} inherits from \texttt{Exception}, and \texttt{Exception} inherits from \texttt{BaseException}, the same object also counts as an \texttt{Exception} and as a \texttt{BaseException}.

The class route can be inspected with \texttt{\_\_mro\_\_}. MRO means method resolution order: the ordered chain of classes Python searches through for inherited behavior.

\begin{verbatim}
print(ValueError.__mro__)
print(KeyError.__mro__)
print(IndexError.__mro__)
\end{verbatim}

Possible output:

\begin{verbatim}
(<class 'ValueError'>, <class 'Exception'>, <class 'BaseException'>, <class 'object'>)
(<class 'KeyError'>, <class 'LookupError'>, <class 'Exception'>, <class 'BaseException'>, <class 'object'>)
(<class 'IndexError'>, <class 'LookupError'>, <class 'Exception'>, <class 'BaseException'>, <class 'object'>)
\end{verbatim}

This tells us that \texttt{KeyError} and \texttt{IndexError} are both more specific forms of \texttt{LookupError}. Therefore, a handler for \texttt{LookupError} can catch both dictionary missing-key failures and sequence invalid-index failures.

\begin{verbatim}
data = ["soup", "pretzels"]

try:
    print(data[42])
except LookupError as err:
    print("lookup failed")
    print(f"type(err): {type(err)}")
    print(f"message: {err}")
\end{verbatim}

Possible output:

\begin{verbatim}
lookup failed
type(err): <class 'IndexError'>
message: list index out of range
\end{verbatim}

The actual exception object is an \texttt{IndexError}, but it is caught by \texttt{LookupError} because \texttt{IndexError} is a child class of \texttt{LookupError}.

The exception object also has attributes and inherited methods:

\begin{verbatim}
try:
    score = int("D'oh!")
except ValueError as err:
    print(f"type(err): {type(err)}")
    print(f"dir(err) sample: {dir(err)[:12]}")
    print(f"err.args: {err.args}")
    print(f"str(err): {str(err)}")
    print(f"repr(err): {repr(err)}")
\end{verbatim}

Possible output:

\begin{verbatim}
type(err): <class 'ValueError'>
dir(err) sample: ['__cause__', '__class__', '__context__', '__delattr__',
'__dict__', '__dir__', '__doc__', '__eq__', '__format__', '__ge__',
'__getattribute__', '__getstate__']
err.args: ("invalid literal for int() with base 10: \"D'oh!\"",)
str(err): invalid literal for int() with base 10: "D'oh!"
repr(err): ValueError('invalid literal for int() with base 10: "D\'oh!"')
\end{verbatim}

The \texttt{args} attribute stores the construction arguments of the exception object. The \texttt{str()} conversion gives the human-readable message. The \texttt{repr()} conversion shows a more explicit object representation.

\subsubsection{Matching Specific Exceptions Before General Exceptions}

Python checks \texttt{except} clauses from top to bottom. The first matching handler is selected. After one handler is selected, later handlers in the same \texttt{try} statement are skipped.

This means specific handlers must come before general handlers.

Correct order:

\begin{verbatim}
text = "forty-two"

try:
    number = int(text)
    print(f"number: {number}")
except ValueError:
    print("specific handler: invalid integer text")
except Exception:
    print("general handler: some other ordinary exception")
\end{verbatim}

Possible output:

\begin{verbatim}
specific handler: invalid integer text
\end{verbatim}

The \texttt{ValueError} handler runs because it appears first and matches the raised exception exactly.

Bad order:

\begin{verbatim}
text = "forty-two"

try:
    number = int(text)
    print(f"number: {number}")
except Exception:
    print("general handler: catches the ValueError first")
except ValueError:
    print("specific handler: unreachable here")
\end{verbatim}

Possible output:

\begin{verbatim}
general handler: catches the ValueError first
\end{verbatim}

The \texttt{ValueError} handler is structurally useless here. The broader \texttt{Exception} handler already caught the object.

The same problem appears with \texttt{LookupError}, \texttt{KeyError}, and \texttt{IndexError}.

\begin{verbatim}
scores = {
    "Jerry": 72,
    "Elaine": 91,
}

name = "Kramer"

try:
    print(scores[name])
except KeyError:
    print("specific handler: missing dictionary key")
except LookupError:
    print("general lookup handler")
\end{verbatim}

Possible output:

\begin{verbatim}
specific handler: missing dictionary key
\end{verbatim}

Here the specific \texttt{KeyError} handler comes first. That is correct if missing dictionary keys need special treatment.

If the broader handler comes first, it catches the same exception:

\begin{verbatim}
scores = {
    "Jerry": 72,
    "Elaine": 91,
}

name = "Kramer"

try:
    print(scores[name])
except LookupError:
    print("general lookup handler")
except KeyError:
    print("specific handler: unreachable here")
\end{verbatim}

Possible output:

\begin{verbatim}
general lookup handler
\end{verbatim}

The practical rule is:

\begin{quote}
Place narrow exception classes first and broad exception classes later.
\end{quote}

This is not a style preference. It follows directly from the handler-selection algorithm.

\subsubsection{Capturing Exception Objects with \texttt{except SomeError as e}}

An \texttt{except} clause can bind the caught exception object to a name. The usual syntax is:

\begin{verbatim}
except SomeError as e:
\end{verbatim}

The name \texttt{e} is not special. It is just a local name bound to the exception object. More descriptive names such as \texttt{err} are often clearer in teaching code.

\begin{verbatim}
values = ["10", "twenty", "42"]

for text in values:
    try:
        number = int(text)
        print(f"{text!r:>8} -> {number:>3}")
    except ValueError as err:
        print(f"{text!r:>8} -> conversion failed")
        print(f"error type: {type(err)}")
        print(f"error args: {err.args}")
\end{verbatim}

Possible output:

\begin{verbatim}
    '10' ->  10
'twenty' -> conversion failed
error type: <class 'ValueError'>
error args: ("invalid literal for int() with base 10: 'twenty'",)
    '42' ->  42
\end{verbatim}

Capturing the object is useful when the program needs the error message, the exact type, the construction arguments, or traceback-related information.

Another example with dictionary lookup:

\begin{verbatim}
scores = {
    "Jerry": 72,
    "Elaine": 91,
    "Kramer": 42,
}

names = ["Jerry", "Newman", "Kramer"]

for name in names:
    try:
        score = scores[name]
        print(f"{name:<8} | {score:>3}")
    except KeyError as err:
        print(f"{name:<8} | missing key object: {err}")
        print(f"err.args: {err.args}")
\end{verbatim}

Possible output:

\begin{verbatim}
Jerry    |  72
Newman   | missing key object: 'Newman'
err.args: ('Newman',)
Kramer   |  42
\end{verbatim}

For a \texttt{KeyError}, the missing key is commonly stored in \texttt{err.args}. In this case, the missing key is the string \texttt{"Newman"}.

The captured exception object exists only because the exceptional route was taken. If no exception occurs, the \texttt{except} block is not entered, and the exception name is not created by that block.

\begin{verbatim}
text = "42"

try:
    number = int(text)
    print(f"number: {number}")
except ValueError as err:
    print(f"conversion failed: {err}")

print("after try")
\end{verbatim}

Possible output:

\begin{verbatim}
number: 42
after try
\end{verbatim}

No \texttt{ValueError} object was created here, because the conversion succeeded.

Capturing is also useful when several exception types share one handler:

\begin{verbatim}
items = ["soup", "pretzels"]

cases = [0, 42, "bad"]

for index in cases:
    try:
        item = items[index]
        print(f"index {index!r}: {item}")
    except (IndexError, TypeError) as err:
        print(f"index {index!r}: lookup failed")
        print(f"error type: {type(err)}")
        print(f"message: {err}")
\end{verbatim}

Possible output:

\begin{verbatim}
index 0: soup
index 42: lookup failed
error type: <class 'IndexError'>
message: list index out of range
index 'bad': lookup failed
error type: <class 'TypeError'>
message: list indices must be integers or slices, not str
\end{verbatim}

The tuple in:

\begin{verbatim}
except (IndexError, TypeError) as err:
\end{verbatim}

means that either exception class can be handled by the same block. The captured object still keeps its exact runtime type.

\subsubsection{The Risk of Bare \texttt{except} and Over-Broad Exception Traps}

A bare \texttt{except} clause has no exception class:

\begin{verbatim}
try:
    risky_operation()
except:
    print("something went wrong")
\end{verbatim}

This catches too much. It can catch exceptions that usually should not be swallowed, including interruption and interpreter-exit signals such as \texttt{KeyboardInterrupt} and \texttt{SystemExit}. For normal application errors, catching \texttt{Exception} is already broad; bare \texttt{except} is broader still.

A safer teaching rule is:

\begin{quote}
Catch the exception class that represents the failure you actually know how to handle.
\end{quote}

Bad pattern:

\begin{verbatim}
texts = ["10", "20", "bad", "42"]

for text in texts:
    try:
        number = int(text)
        print(f"{text!r} -> {number}")
    except:
        print("ignored an error")
\end{verbatim}

Possible output:

\begin{verbatim}
'10' -> 10
'20' -> 20
ignored an error
'42' -> 42
\end{verbatim}

The program hides what kind of error occurred. It also makes future bugs harder to diagnose, because unrelated mistakes would be swallowed by the same handler.

Better pattern:

\begin{verbatim}
texts = ["10", "20", "bad", "42"]

for text in texts:
    try:
        number = int(text)
        print(f"{text!r} -> {number}")
    except ValueError as err:
        print(f"{text!r} -> invalid integer text")
        print(f"message: {err}")
\end{verbatim}

Possible output:

\begin{verbatim}
'10' -> 10
'20' -> 20
'bad' -> invalid integer text
message: invalid literal for int() with base 10: 'bad'
'42' -> 42
\end{verbatim}

Now the handler says exactly what it expects: \texttt{int()} may fail because the text is not a valid integer.

Over-broad \texttt{Exception} handlers can also hide bugs:

\begin{verbatim}
items = ["10", "20", "42"]

try:
    for item in items:
        number = int(item)
        total = total + number
except Exception as err:
    print("some error happened")
    print(f"type(err): {type(err)}")
    print(f"message: {err}")
\end{verbatim}

Possible output:

\begin{verbatim}
some error happened
type(err): <class 'NameError'>
message: name 'total' is not defined
\end{verbatim}

The real problem is not bad input. The real problem is a programmer mistake: \texttt{total} was used before being assigned. A broad handler catches it, but catching it does not make the program correct.

A better version initializes the state and catches only conversion failures:

\begin{verbatim}
items = ["10", "20", "42"]
total = 0

for item in items:
    try:
        number = int(item)
    except ValueError as err:
        print(f"invalid integer text: {item!r}")
        print(f"message: {err}")
    else:
        total = total + number

print(f"total: {total}")
\end{verbatim}

Possible output:

\begin{verbatim}
total: 72
\end{verbatim}

The handler is now narrow. It catches the failure the program intends to recover from. Other programming mistakes are allowed to surface.

There are cases where a broad handler is useful, such as logging an unexpected failure before re-raising it. But swallowing a broad exception silently is usually dangerous.

Acceptable diagnostic pattern:

\begin{verbatim}
try:
    number = int("D'oh!")
except Exception as err:
    print("unexpected failure during conversion block")
    print(f"type(err): {type(err)}")
    print(f"message: {err}")
    raise
\end{verbatim}

Possible output:

\begin{verbatim}
unexpected failure during conversion block
type(err): <class 'ValueError'>
message: invalid literal for int() with base 10: "D'oh!"
Traceback (most recent call last):
  ...
ValueError: invalid literal for int() with base 10: "D'oh!"
\end{verbatim}

The handler records information, then \texttt{raise} sends the same exception onward. Re-raising will be discussed later in the propagation section.

The practical rules are:

\begin{itemize}
    \item Avoid bare \texttt{except} in ordinary application code.
    \item Avoid broad \texttt{except Exception} unless there is a clear reason.
    \item Catch specific exception classes when recovery is specific.
    \item Do not use exception handling to hide programming mistakes.
    \item If a broad handler is used for logging, usually re-raise the exception afterward.
\end{itemize}

For control-flow graph analysis, a bare or over-broad handler creates a very large exceptional edge. Many unrelated failures are routed into the same block. That may keep the program running, but it can also destroy the information needed to understand why the program failed.